\subsubsection{P3：时空网络、容量约束与资源受限调度}

P3 对应的核心运筹学问题，是在已给定任务分配和任务顺序后，进一步判断多台 AMR 在时间和空间上的资源占用是否可行。P2 更偏向任务层面的整数规划，P3 则把计划放入“时间--空间--资源”框架中检查，重点体现时空网络建模、容量约束、互斥约束、资源受限调度和启发式修复等知识点。

\paragraph{时空网络建模。}
时间扩展网络是运筹学中处理动态路径和动态占用问题的常用方法。它把普通网络中的节点和边复制到不同时间层，使“车辆在什么时刻位于什么位置”能够被显式表示。P3 将 P2 的任务级计划展开为带绝对时间的二维轨迹，本质上就是构造离散时间下的时空网络占用表。这样，路径不再只是起点到终点的几何连接，而是变成 AMR 在各个时刻对仓库空间资源的连续占用。

\paragraph{容量约束。}
容量约束用于限制同一资源在同一时刻能够被多少个对象占用。在 P3 中，取货点、分拣点、出库点等关键节点可视为网络资源，其容量通常为 1，可写成
\[
  \sum_{k\in K} y_{kvt}\le c_v .
\]
其中 $y_{kvt}$ 表示 AMR $k$ 在时刻 $t$ 是否占用节点 $v$，$c_v$ 表示节点容量。P3 对 \verb|P*|、\verb|SORT*|、\verb|OUT*| 等节点的占用检测，就是这一类网络容量约束在仓储调度中的体现。

\paragraph{互斥约束与安全距离约束。}
多 AMR 运行还需要满足空间互斥条件，即两台机器人不能在同一时间占用过近的空间。P3 用 footprint 半径和安全裕量构造距离约束：
\[
  \|X_i(t)-X_j(t)\|\ge \rho_i+\rho_j+\sigma ,
  \quad i\ne j .
\]
这属于运筹学中常见的冲突避免约束或互斥资源约束。与只限制图上同一节点、同一边不同，P3 直接在二维坐标上判断 AMR 间距，因此能够反映仓库通道宽度、货架间距和机器人 footprint 对可执行性的影响。

\paragraph{资源受限调度。}
从调度理论看，P3 可理解为资源受限调度问题。AMR 是执行任务的机器资源，仓库节点和通道是共享空间资源，空驶、取货、载货和等待都是时间轴上的作业活动。P2 给出的任务顺序只是初始作业序列，P3 需要检查这些作业活动叠加后是否超出共享资源容量。若发生冲突，就说明任务级计划虽然可行，但在资源受限的执行层面仍不可行。

\paragraph{启发式算法与局部搜索。}
P3 没有重新求解全局任务分配模型，而是在 P2 计划附近进行局部修复。等待、换路和同车相邻换序分别对应时间调整、路径邻域搜索和序列邻域搜索。算法每轮选择当前冲突，生成有限候选动作，再比较修复后的冲突数和目标值。这体现了启发式算法在大规模组合调度中的思想：不追求一次性建立过大的全局精确模型，而是利用局部搜索快速得到可执行方案。

\paragraph{目标规划与字典序优化。}
P3 的评价规则体现目标规划思想。剩余冲突数具有最高优先级，因为存在冲突的方案不能执行；在冲突数相同的情况下，再比较延期、最大完工时间 $C_{\max}$、新增等待和对原计划的扰动。这样的字典序结构避免了用较短完工时间去交换安全冲突，也避免为了消除局部冲突而过度破坏 P2 的任务安排。最终，P3 将任务级计划转化为满足容量约束和安全互斥约束的轨迹级可执行计划。
